\fancyhead[L]{\small ANZEIGEN.}
\SourceLeaf{346}
\pdfbookmark[0]{Anzeigen. Disquisitiones generales circa superficies curvas.}{anzeigen-curvas}
\begin{center}
\Large ANZEIGEN.\par\medskip
\normalsize\rule{.85\textwidth}{.3pt}\par
Göttingische gelehrte Anzeigen.\quad 1827 November 5.\par
\rule{.85\textwidth}{.3pt}
\end{center}
\hspace*{\parindent}Am 8. October überreichte Hr. Hofr. \textsc{Gauss} der Königl. Societät eine Vorlesung:
\begin{center}\textit{Disquisitiones generales circa superficies curvas.}\end{center}
\hspace*{\parindent}Obgleich die Geometer sich viel mit allgemeinen Untersuchungen über die krummen Flächen beschäftigt haben, und ihre Resultate einen bedeutenden Theil des Gebiets der höhern Geometrie ausmachen, so ist doch dieser Gegenstand noch so weit davon entfernt, erschöpft zu sein, dass man vielmehr behaupten kann, es sei bisher nur erst ein kleiner Theil eines höchst fruchtbaren Feldes angebauet. Der Verf. hat schon vor einigen Jahren durch die Auflösung der Aufgabe, alle Darstellungen einer gegebenen Fläche auf einer andern zu finden, bei welchen die kleinsten Theile ähnlich bleiben, dieser Lehre eine neue Seite abzugewinnen gesucht: der Zweck der gegenwärtigen Abhandlung ist, abermals andere neue Gesichtspunkte zu eröffnen, und einen Theil der neuen Wahrheiten, die dadurch zugänglich werden, zu entwickeln. Wir werden davon hier anzeigen, was ohne zu grosse Weitläuftigkeit verständlich gemacht werden kann, müssen aber dabei im Voraus bemerken, dass sowohl die neuen Begriffsbildungen, als die Theoreme, wenn die grösste Allgemeinheit umfasst werden soll, zum Theil noch einiger Beschränkungen oder nähern Bestimmungen bedürfen, welche hier übergangen werden müssen.
